\begin{theorem}[Second main term and witness sets from an isolated dominant channel]\label{thm:kernel-chebotarev-second-main-term-witness}
Fix a finite group \(G\) and a temperature parameter \(u\in(0,1]\) in the setting of Assumption \ref{thm:kernel-chebotarev-exp}.
Let \(\lambda_1(u)=\spr(M_{K,\ind}(u))\) be the Perron root in the trivial channel, and for each irreducible representation \(\rho\neq \ind\) write \(M_{K,\rho}(u)\) for the corresponding twisted channel matrix.
Assume that there exists a nontrivial channel \(\rho_\star\neq \ind\) and a complex number \(\mu_\star\in\mathrm{spec}(M_{K,\rho_\star}(u))\) such that
\begin{enumerate}
  \item \(\abs{\mu_\star}=\eta(u)>0\);
  \item across all nontrivial channels, the spectral values of modulus \(\eta(u)\) consist exactly of the pair \(\{\mu_\star,\overline{\mu_\star}\}\), counted with algebraic multiplicity, where the two values may coincide if \(\mu_\star\in\RR\);
  \item each dominant spectral value in \(\{\mu_\star,\overline{\mu_\star}\}\) is simple; and
  \item
  \[
  \Gamma(u):=\max\bigl\{\abs{\nu}:\ \nu\in\mathrm{spec}(M_{K,\rho}(u)),\ \rho\neq \ind,\ \nu\notin\{\mu_\star,\overline{\mu_\star}\}\bigr\}
  \]
  satisfies \(\Gamma(u)<\eta(u)\).
\end{enumerate}
Then for every conjugacy-invariant subset \(A\subseteq G\), with
\[
B_{n,A}(u):=\sum_{c\in A}B_{n,c}(u),
\]
there exists a complex coefficient \(c_A\) such that
\[
B_{n,A}(u)
=
\mathrm{unif}_G(A)\frac{\lambda_1(u)^n}{n}
\;+\;
\Re\!\bigl(c_A\,\mu_\star^{\,n}\bigr)\frac{1}{n}
\;+\;
O\!\left(\frac{\Gamma(u)^n}{n}\right).
\]
If, in addition, \(c_A\neq 0\), then there exists a constant \(c(A)>0\) such that
\[
\bigl|\Re(c_A\mu_\star^{\,n})\bigr|\ \ge\ c(A)\,\eta(u)^n
\]
for infinitely many integers \(n\).
\end{theorem}

\begin{proof}
For each conjugacy-invariant subset \(A\subseteq G\), Assumption \ref{thm:kernel-chebotarev-exp} provides linear functionals \(L_{A,\rho}\) such that
\[
B_{n,A}(u)
=
\mathrm{unif}_G(A)\frac{\lambda_1(u)^n}{n}
\;+\;
\frac{1}{n}\sum_{\rho\in\widehat G\setminus\{\ind\}}L_{A,\rho}\!\bigl(M_{K,\rho}(u)^n\bigr).
\]
Fix \(\rho\neq \ind\), and write \(T_\rho:=M_{K,\rho}(u)\). By the Jordan decomposition of the finite-dimensional operator \(T_\rho\), the spectral assumptions imply that
\[
T_\rho^n
=
\sum_{\omega\in\mathrm{spec}(T_\rho)\cap\{\mu_\star,\overline{\mu_\star}\}}\omega^n P_{\rho,\omega}
+R_{\rho,n},
\]
where \(P_{\rho,\omega}\) is the spectral projector onto the one-dimensional eigenspace for \(\omega\), and \(\|R_{\rho,n}\|=O(\Gamma(u)^n)\) with respect to any fixed matrix norm. Since there are only finitely many nontrivial channels and the functionals \(L_{A,\rho}\) are fixed, summing over \(\rho\neq \ind\) gives coefficients \(a_{A,\omega}\in\CC\) such that
\[
\sum_{\rho\in\widehat G\setminus\{\ind\}}L_{A,\rho}\!\bigl(M_{K,\rho}(u)^n\bigr)
=
\sum_{\omega\in\{\mu_\star,\overline{\mu_\star}\}}a_{A,\omega}\,\omega^n
+O(\Gamma(u)^n).
\]
If \(\mu_\star\notin\RR\), then the left-hand side is real for every \(n\), so necessarily
\[
a_{A,\overline{\mu_\star}}=\overline{a_{A,\mu_\star}}.
\]
Setting \(c_A:=2a_{A,\mu_\star}\) therefore yields
\[
\sum_{\rho\in\widehat G\setminus\{\ind\}}L_{A,\rho}\!\bigl(M_{K,\rho}(u)^n\bigr)
=
\Re\!\bigl(c_A\mu_\star^n\bigr)+O(\Gamma(u)^n),
\]
whereas if \(\mu_\star\in\RR\), then the dominant set consists only of \(\mu_\star\), and the reality of the left-hand side shows that the corresponding coefficient \(a_{A,\mu_\star}\) is real. In that case the same formula holds with \(c_A:=a_{A,\mu_\star}\in\RR\). Substituting into the preceding display proves the second-main-term expansion.

Finally, \(\mu_\star/\eta(u)\in\TT\), so its powers either lie in a finite subgroup of the circle or are dense in one. In either case there are infinitely many integers \(n\) for which \(\abs{\Re(c_A\mu_\star^n)}\) is bounded below by a positive constant times \(\eta(u)^n\), which gives the witness lower bound.
\end{proof}

\endinput
